% \usepackage[colorinlistoftodos,prependcaption,textsize=tiny]{todonotes}
% \usepackage{xargs} 

\usepackage{wrapfig}
\usepackage{threeparttable}
\usepackage{enumitem}


% Command for the name of the proposed method
\newcommand{\our}{\texttt{ReEvo}}
% \newcommand{\our}{\textsc{ReEvo}}
% you may also simpy have no sty.e 


% Note: you can add your own (these are the TODO notes)
% \newcommandx{\FB}[2][1=]{\todo[linecolor=orange,backgroundcolor=orange!25,bordercolor=orange,#1]{FB: #2}}

% \newcommandx{\HY}[2][1=]{\todo[linecolor=blue,backgroundcolor=blue!25,bordercolor=blue,#1]{HY: #2}}

\newcommand{\fede}[1]{\textcolor{orange}{#1}}


%%% Prompts
% Define additional custom colors
\definecolor{promptbackground}{rgb}{0.98, 0.98, 0.98} % Lighter background for prompts
\definecolor{textcolor}{rgb}{0, 0, 0} % Darker text color for better readability
\definecolor{backcolour}{rgb}{0.95, 0.95, 0.92}
\definecolor{codegreen}{rgb}{0,0.6,0}
\definecolor{codegray}{rgb}{0.5,0.5,0.5}
\definecolor{codepurple}{rgb}{0.58,0,0.82}
\definecolor{codemagenta}{rgb}{0.78,0,0.52}
\definecolor{codeblue}{rgb}{0.1,0.1,0.9}

% Custom style for prompts
\lstdefinestyle{promptstyle}{
    backgroundcolor=\color{promptbackground},
    basicstyle=\tiny\color{textcolor}, % Italic and footnotesize for text snippets
    breakatwhitespace=true,         % Only break lines at whitespaces
    breaklines=true,                % Enable line breaking
    captionpos=b,                   % Position caption at the bottom
    keepspaces=true,                % Preserve spaces for alignment
    showspaces=false,               % Don't show spaces explicitly
    showstringspaces=false,         % Don't show spaces in strings
    showtabs=false,                 % Don't show tabs
    frame=single,                   % Single frame around the text
    rulecolor=\color{textcolor},    % Use text color for frame
    framesep=3pt,                   % Padding between text and frame
    frameround=tttt,                % All corners rounded
    framexleftmargin=5pt,           % Extra margin on the left for symmetry
    xleftmargin=5pt,                % Left margin for alignment
    xrightmargin=5pt,               % Right margin for alignment
    tabsize=2,                      % Standard tab size
    linewidth=\textwidth            % Line width fit to page
}



%%% Prompt code
% Define custom colors

% Custom style. Comment out for default
\lstdefinestyle{promptcodestyle}{
    backgroundcolor=\color{promptbackground},
    commentstyle=\color{codegreen},
    keywordstyle=\color{codemagenta}, % Changed from magenta to custom magenta
    numberstyle=\tiny\color{codegray},
    stringstyle=\color{codepurple},
    basicstyle=\tiny\ttfamily, % Changed from tiny to footnotesize for better readability
    breakatwhitespace=false,
    breaklines=true,
    captionpos=b,
    keepspaces=true,
    numbersep=8pt, % Increased separation between numbers and text for clarity
    showspaces=false,
    showstringspaces=false,
    showtabs=false,
    tabsize=2, % Increased tab size for better alignment
    frame=single, % Adds a frame around the code
    frameround=tttt,                % All corners rounded
    rulecolor=\color{black}, % Color of the frame
}



%%% Generated heuristics

% Custom style. Comment out for default
\lstdefinestyle{heuristicstyle}{
    backgroundcolor=\color{backcolour},
    commentstyle=\color{codegreen},
    keywordstyle=\color{codemagenta}, % Changed from magenta to custom magenta
    numberstyle=\tiny\color{codegray},
    stringstyle=\color{codepurple},
    basicstyle=\tiny\ttfamily, % Changed from tiny to footnotesize for better readability
    breakatwhitespace=false,
    breaklines=true,
    captionpos=b,
    keepspaces=true,
    numbersep=8pt, % Increased separation between numbers and text for clarity
    showspaces=false,
    showstringspaces=false,
    showtabs=false,
    tabsize=2, % Increased tab size for better alignment
    frame=single, % Adds a frame around the code
    rulecolor=\color{black}, % Color of the frame
}

% \DeclareMathOperator{\tanh}{tanh} % Define \tanh as a math operator

\newcommand{\red}[1]{\textcolor{red}{#1}}